\begin{statement}{C.2.UN.48}{Proved}{The full compatible cycle is realized over a compact rational base}
Take the specified cycle of primary or secondary diagonal
Hopf surfaces of algebraic dimension zero in C.2.UN.46,
with its specified elliptic gluing isomorphisms.
Suppose its dual-normal bundles are isomorphic to the
trivial bundle, and compute the integral edge product
from those actual isomorphisms as in that statement.
Every cycle whose product is the identity occurs,
with precisely its given components and gluing maps,
as a reduced SNC fibre of a holomorphic map
\[
                 f:Y\longrightarrow\PP^1
\]
from a smooth compact complex threefold with connected
two-torus general fibre. Every reduced fibre of this
map has Euler characteristic zero.

In particular the necessary criterion of C.2.UN.46
is sufficient for compact occurrence of the whole
reduced support, for arbitrary cycle length and for
secondary components as well. The construction makes
specified choices of the normal-product constants;
it preserves the entire reduced surface, while
recording those choices and the actual punctured-family
comparison below.

For example there is a three-component compact
realization with primary Hopf multipliers
\[
 \begin{gathered}
 (e^{-1},e^{-\sqrt2}),\qquad
 (e^{-\sqrt2},e^{-1-\sqrt2}),\qquad
 (e^{-1-\sqrt2},e^{-1}),\\
 L_1=\begin{pmatrix}0&-1\\-1&-1\end{pmatrix},\qquad
 L_2=\begin{pmatrix}0&1\\-1&1\end{pmatrix},\qquad
 L_3=\begin{pmatrix}0&-1\\-1&1\end{pmatrix}.
 \end{gathered}
 \tag{C.HG.1}
\]
Here $L_3L_2L_1=I$ and the actual monodromy period
on the common open group is $(t^{-2},t^2)$.
Thus the longer-cycle occurrence assertion includes
an explicit case beyond the two-component inversion
construction.
\end{statement}

\begin{proof}
\par\medskip\noindent\textbf{The common lattice and the sign of the normal periods.}
Retain the exact logarithmic coordinates of C.2.UN.46.
The product identity permits constant matrices $P_i$
identifying every open group with one common group
\[
                  G=\C^2/\Lambda,\qquad \Lambda\simeq\Z^3,
                  \qquad P_{i+1}L_i=P_i.
\tag{C.HG.2}
\]
Take indices cyclically. Write $n_i=P_i(2\pi i,0)$
for the outgoing normal meridian. Since
$L_i(1,0)=(0,-1)$, the incoming meridian on the
$i$th component is $-n_{i-1}$, and
\[
 P_i=\left(\frac{n_i}{2\pi i},
                   -\frac{n_{i-1}}{2\pi i}\right),
 \qquad
 2\pi i v_0=-\sum_i n_i=:u\in\Lambda.
\tag{C.HG.3}
\]
The last equality follows by adding the actual
edge translations $P_{i+1}(d_i,e_i+\log t)$.
It retains their signs, rather than choosing a
meridian sign afterwards.

Put $W=\operatorname{span}_{\R}\Lambda$, and choose
a nonzero real linear functional $\ell:\C^2\to\R$
with kernel $W$. On the $i$th Hopf surface choose
its actual contracting generator with multipliers
$\alpha_i,\beta_i$, including when a finite diagonal
subgroup is also present. Its common logarithm
\[
 c_i=P_i(\log\alpha_i,\log\beta_i)
       =\frac{\log\alpha_i}{2\pi i}n_i
          -\frac{\log\beta_i}{2\pi i}n_{i-1}
\]
belongs to $\Lambda$. Applying $\ell$ and using
$\ell(n_i)=\ell(n_{i-1})=0$ gives the exact real
identity
\[
 \log|\alpha_i|\,\ell(i n_i)
       =\log|\beta_i|\,\ell(i n_{i-1}).
\tag{C.HG.4}
\]
Both logarithms are strictly negative.
Neither displayed value of $\ell$ is zero: if
both vanished, $\ell$ would vanish on the complex
span of $n_i,n_{i-1}$, which is $\C^2$ by
(C.HG.3). The identity would make both zero if
one were zero. It follows successively around
the cycle that all $\ell(i n_i)$ have the same
nonzero sign. Consequently
\[
                       \ell(iu)=-\sum_i\ell(i n_i)\ne0.
\tag{C.HG.5}
\]
In particular $u\ne0$, and $\C u\cap W=\R u$.

Write $u=d p$, where $p$ is a primitive element
of the free lattice $\Lambda$ and $d>0$.
The real quotient map
$W/\R p\to\C^2/\C p$ is an isomorphism, by
(C.HG.5). Therefore the image of $\Lambda$ in
$\C^2/\C p$ is a discrete rank-two lattice.
Its quotient is an actual elliptic curve.
Extend $p$ to an integral basis $p,\lambda,c$
of $\Lambda$. Choose a linear quotient
coordinate $X$ with $X(\lambda)=2\pi i$.
The two projected periods are real-linearly
independent, so $\Re X(c)\ne0$; reverse $c$
if necessary to make it negative. Choose a
second complex linear coordinate $Z$ with
$Z(p)=2\pi i$ and $Z(\lambda)=0$. These
conditions determine it because $p,\lambda$
are complex-linearly independent. In these
actual linear coordinates the common lattice
is
\[
 \Lambda=\left\langle
     (0,2\pi i),\ (2\pi i,0),\
     (\log q_0,\log a)\right\rangle_{\Z},
 \qquad 0<|q_0|<1,\qquad a\in\C^*,
 \qquad u=d(0,2\pi i).
\tag{C.HG.6}
\]
The logarithms in this formula are the actual
chosen coordinates of $c$.

Here $a\ne1$. Indeed, if $a=1$, subtract an
integral multiple of the first period from $c$
so that its second coordinate is zero. The map
\[
 \C/\langle2\pi i,\log q_0\rangle_{\Z}
          \longrightarrow G,\qquad [X]\longmapsto[(X,0)]
\tag{C.HG.7}
\]
is well defined and injective: equality in $G$
forces the coefficient of $(0,2\pi i)$ to vanish,
leaving exactly the two periods in its source.
It is an immersion and, by compactness, an embedded
elliptic curve. But $G$ is the complement of the
two axes in any of the prescribed Hopf surfaces.
The curve calculation of C.2.UN.40 and C.2.UN.43
says that these are all its compact curves.
The image of (C.HG.7) contradicts that calculation.
This proves the precise specialization needed
for the pencil construction below.

\par\medskip\noindent\textbf{A compact elliptic suspension with the required limiting periods.}
The cubic construction in C.2.UN.41 works for
every specified $a\in\C^*\setminus\{1\}$:
its choice of the smooth nodal point uses no
inequality on $|a|$. Explicitly use
\[
 \begin{gathered}
 y^2z=x^2(x+z),\qquad O=[0:1:0],\\
 P_0=\left(\frac{4a}{(1-a)^2},
                  \frac{4a(1+a)}{(1-a)^3}\right)
 \end{gathered}
\tag{C.HG.8}
\]
in the affine chart $z=1$, whose smooth group
coordinate is $a$. The point is distinct from
$O$ and from the node. Choose a second cubic
through $O,P_0$, transverse at nine distinct
smooth points of the nodal cubic and missing
its node. The jet-surjectivity and nonempty
open conditions in C.2.UN.41 give this choice.
Blowing up the nine actual base points gives
a compact rational elliptic surface with
sections $O,P$ and the selected $I_1$ fibre.

Make the global base change $\PP^1_s\to\PP^1$
which is $s\mapsto s^d$ at that value.
Normalize, resolve and take the relatively
minimal elliptic model $J\to\PP^1_s$.
Near the chosen node the equation before
resolution is $xy=s^d$. Its resolution has
charts $(u_j,v_j)$, $0\le j<d$, with exact maps
\[
 x=u_j^{j+1}v_j^j,\qquad
 y=u_j^{d-j-1}v_j^{d-j},\qquad s=u_jv_j.
\]
Adjacent charts glue by
$u_{j+1}=v_j^{-1}$, $v_{j+1}=u_jv_j^2$.
These maps are isomorphisms off the node and
insert the chain of $d-1$ projective lines
between its two branches; each displayed
central equation is $u_jv_j=0$.
Together with the old component
this gives an $I_d$ fibre. The sections avoid
these charts and retain their given values.
At all other values these operations are
proper vertical modifications of a compact
elliptic surface. The section translation by
$P$ extends to an automorphism of this
relatively minimal model, by the exact graph
and relative-canonical-nef argument of
C.2.UN.41 and C.2.UN.43.

The analytic Tate comparison at this fibre
gives, on a sufficiently small punctured disc,
\[
       J_s=\C^*/Q(s)^{\Z},\qquad
       \operatorname{ord}_0Q=d,\qquad P_s=[a(s)],
       \qquad a(0)=a.
\tag{C.HG.9}
\]
The lift $a(s)$ is a holomorphic unit, since
the pulled-back section meets the smooth
identity component. Choose the holomorphic
$d$th root of the unit $Q(s)/s^d$ and use the
actual coordinate
$t=s\bigl(Q(s)/s^d\bigr)^{1/d}$.
It has nonzero derivative at zero and makes
$Q=t^d$ exactly. Denote the section lift in
this coordinate by $a(t)$, and fix the
holomorphic logarithm with value $\log a$.

Now take the actual suspension
\[
 Y^{\rm std}=
  (\C^*_b\times J)/
      \langle(b,y)\longmapsto(q_0 b,y+P)\rangle
          \longrightarrow\PP^1_s .
\tag{C.HG.10}
\]
The action is free and properly discontinuous
because its first coordinate has modulus
multiplier $|q_0|<1$. A closed annulus in
$\C^*$ times compact $J$ meets every orbit,
so the quotient is compact and smooth.
Its general connected fibres are two-tori.
Near the distinguished value their exact
logarithmic periods are
\[
 \Lambda(t)+\Z(0,d\log t),\qquad
 \Lambda(t)=
 \left\langle(0,2\pi i),(2\pi i,0),
                  (\log q_0,\log a(t))\right\rangle_{\Z}.
\tag{C.HG.11}
\]
This is the displayed suspension quotient,
not just equality of characteristic polynomials.

\par\medskip\noindent\textbf{Transporting the prescribed Hopf surfaces and the whole cycle.}
View each original $n_i$ and contracting
$c_i$ as its actual integral combination
of the three basis periods in (C.HG.6).
Use those same integer coefficients in
$\Lambda(t)$ to define $n_i(t),c_i(t)$.
All integral relations are retained exactly.
In particular $\sum_i n_i(t)=-d(0,2\pi i)$.
Set
\[
 \begin{gathered}
 P_i(t)=\left(\frac{n_i(t)}{2\pi i},
                    -\frac{n_{i-1}(t)}{2\pi i}\right),\\
 \Lambda_i(t)=P_i(t)^{-1}\Lambda(t),
 \quad L_i(t)=P_{i+1}(t)^{-1}P_i(t).
 \end{gathered}
\tag{C.HG.12}
\]
These matrices are holomorphic and invertible
after shrinking the disc. The first column
of $L_i(t)$ is $(0,-1)$, so it has exactly the
form
$\left(\begin{smallmatrix}0&c_i^*(t)\\-1&k_i^*(t)\end{smallmatrix}\right)$
required for the normal-product map.
Its cycle product is the identity.

The group $\Lambda_i(t)/(2\pi i\Z^2)$ is
the same abstract finitely generated group
as at zero, by the fixed integral relations.
Its selected free generator is the actual
diagonal contraction with logarithms
$P_i(t)^{-1}c_i(t)$. Both moduli remain
strictly between zero and one on a smaller
disc. Every finite torsion generator has
a relation $n v=A(2\pi i,0)+B(0,2\pi i)$
with the same integers as at zero. Therefore
its two roots of unity and their coprimality
conditions remain unchanged. These are
actual holomorphic families of the prescribed
primary or secondary Hopf quotients. The
contractions are uniformly contracting over
a closed smaller disc, and their fixed
finite groups remain free. Thus the
quotients of $(\C^2\setminus\{0\})\times\Delta$
are smooth proper families, with precisely
the original surfaces at zero.

Let $d_i$ be the constants specifying the
original elliptic gluing maps in the
coordinates of C.2.UN.46. Keep these constants.
Their normal multipliers have freely
adjustable nonzero constants, whose logarithms
we denote by $e_i(t)$. The full affine cycle
translation in common coordinates is
\[
 \sum_iP_{i+1}(t)
        \binom{d_i}{e_i(t)+\log t}
       =b_0(t)+(0,d\log t).
\tag{C.HG.13}
\]
Two adjacent coefficient vectors of the
unknowns $e_i(t)$ are
$-n_i(t)/(2\pi i)$ and $-n_{i+1}(t)/(2\pi i)$.
They are complex-linearly independent by
(C.HG.12). Fix all the other $e_i(t)$, and
solve the resulting invertible two-by-two
linear system holomorphically to obtain
\[
                              b_0(t)=0.
\tag{C.HG.14}
\]
Only the normal-product constants change.
The central surfaces and every given
elliptic gluing map remain exactly the
prescribed ones.

Use the actual varying normal line bundles
and these affine edge maps to perform the
finite plumbing of C.2.UN.46, with $xy=t$
in each neck. The flat metrics, disjoint
axis neighbourhoods and compact cores may
be chosen uniformly on a smaller parameter
disc; the strict inequalities used in that
proof persist by continuity. The result
$\mathcal Y\to\Delta_t$ is a smooth proper
threefold whose reduced central fibre is
the given cycle, with every component
of multiplicity one.

There is an actual punctured-family
biholomorphism to (C.HG.10). To see the map
without suppressing affine data, choose one
core as origin in the common logarithmic
coordinates. On consecutive cores use
$P_i(t)$ and subtract the accumulated edge
translations from (C.HG.13). Use the same
formulas on the necks, where $xy=t$.
They agree on successive overlaps. On the
last overlap their difference is the
period $(0,d\log t)$ by (C.HG.14);
changing logarithm branches changes them
by $\Lambda(t)$. Hence they define a
holomorphic map to the quotient with
periods (C.HG.11). Fibrewise this is the
closed-coframe comparison of C.2.UN.46,
which is a biholomorphism: its periods
are all the generators of the actual
fundamental group. Its differential is
invertible in the fibre directions and
it is the identity on the base, so it
is a local biholomorphism of total
spaces as well. The fibrewise inverses
therefore give the actual holomorphic
inverse over the punctured disc.

\par\medskip\noindent\textbf{Compact replacement and Euler numbers.}
Replace the distinguished tube in
$Y^{\rm std}$ by $\mathcal Y$, using this
biholomorphism on a punctured annulus.
More precisely remove a closed smaller
disc from the base of the first family
and glue the second over a larger disc.
Their annular overlap is the displayed
isomorphism. Properness of both families
and two nested radii give a Hausdorff
gluing: at either boundary of the annulus
only one of the two pieces contains the
limiting base point. Closed smaller
pieces in these two proper families
cover the resulting space, proving
compactness. Smoothness and the
holomorphic map to the same $\PP^1_s$
are preserved. Its fibres are connected
and its general fibre is the specified
suspension two-torus.

The new central support has Euler number
zero, since every Hopf component and every
double elliptic curve has Euler number
zero and there are no triple points.
For every other base value the support
is unchanged from (C.HG.10). Its map to
the actual elliptic curve $\C^*/q_0^{\Z}$
is a locally trivial bundle with fibre
the reduced curve $J_s^{\rm red}$ and
transition the automorphism induced by
$P$. A finite triangulation, or the
Euler characteristic of its finite
local-system cellular complex, gives
Euler number
$e(\C^*/q_0^{\Z})e(J_s^{\rm red})=0$.
Thus every reduced fibre has Euler
number zero, as asserted.

Finally (C.HG.1) is verified by taking,
in the common angular lattice,
\[
 n_1=2\pi i(1,0),\quad
 n_2=2\pi i(1,-1),\quad
 n_3=2\pi i(0,-1),\quad
 c=(-1,-\sqrt2).
\tag{C.HG.15}
\]
The three chosen contracting logarithms
are $P_1^{-1}c$, $-P_2^{-1}c$ and
$-P_3^{-1}c$, exactly the pairs in
(C.HG.1). Their ratios are irrational,
so the three Hopf surfaces have
algebraic dimension zero. The edge
maps in multiplicative coordinates,
with constant factors one, are
\[
 \begin{gathered}
 (z,w)\mapsto(w^{-1},t z^{-1}w^{-1}),\\
 (z,w)\mapsto(w,t z^{-1}w),\\
 (z,w)\mapsto(w^{-1},t z^{-1}w).
 \end{gathered}
\tag{C.HG.16}
\]
Their composite is $(z,w)\mapsto(t^{-2}z,t^2w)$.
The exact integral matrices, negative
contraction logarithms, primitive
period divisor $d=2$ and the two-constant
translation solution are checked in
\path{code/D_hopf_cycle_matrices.py}
and its saved output.

The finite-secondary divisibility is also
explicit. For a repeated two-component
cycle of length $2k$ with finite weights
$(a,b)$ modulo $n$, the common angular
lattice is
$\Z^2+\Z(a,b)/n$, and the period vector is
$k(-1,1)$ in units of $2\pi i$.
Its intersection with $\R(-1,1)$ is
$\Z(-1,1)/\gcd(n,a+b)$: the sum of the
coordinates requires $n\mid j(a+b)$ for
a representative $j(a,b)/n$, and the
coprimality of $a$ to $n$ then gives
exactly those fractions on this line.
Consequently the precise positive divisor is
\[
                    d=k\gcd(n,a+b).
\tag{C.HG.17}
\]
The sign is fixed by the actual outgoing
vectors $(1,0)$ and $(0,-1)$.
For example weights $(1,4)$ modulo $5$
give $d=5$ for two components.
In the three-component example, the
finite weights $(1,2)$ modulo $5$ become
$(1,2),(3,2),(2,4)$ on its three actual
Hopf quotients. They are all free, the
common period still has divisor $2$,
and the same construction gives this
compact secondary three-cycle.
The saved computation checks 12,180
repeated-cycle cases and 1,376 free
secondary three-cycles by exact
Hermite normal forms.
\end{proof}


% C.2.UN.46 from part_C2_hopf_cycles.tex

